%--------------------------------------------------------------------------
\chapter{\label{Conclusion}Conclusion (WORK) (READ)}
%--------------------------------------------------------------------------

%--------------------------------------------------------------------------
\section{\label{Overview}Overview and Contributions}
%--------------------------------------------------------------------------

%--------------------------------------------------------------------------
\subsection{\label{Overview:SED} 
Molecular Dynamics-based Methods for Predicting Vibrational 
Lifetimes}
%--------------------------------------------------------------------------

In Chapter \ref{Chapter:SED}, two MD-based methods for predicting 
phonon properties and thermal conductivity were compared. 
The $\Phi$ method, which is the NMD method in the frequency-domain, 
was properly derived the starting with anharmonic lattice 
dynamics theory (see Appendix \ref{Appendix_A:Derivation}). 
The meaning of the proposed spectral method, $\Phi'$, 
was clarified and related it to the dynamic structure factor 
(see Appendix \ref{Appendix_B}). 
While 
the $\Phi'$ method does not accurately predict the mode lifetimes, 
the advantage of the $\Phi'$ versus the $\Phi$ method is that it does 
not require an eigenvalue solution for the mode eigenvectors.  
The dynamic structure factor, closely related to the $\Phi'$ method 
(see Appendix \ref{Appendix_B}), can predict frequency-dependent 
timescales from MD simulations for systemswith a larger number of 
atoms than those studied in this work using the NMD method 
(see Section \ref{Future:Timescales}). 

%--------------------------------------------------------------------------
\subsection{\label{Overview:VC}
Thermal Transport in Alloys and the High-scatter Limit}
%--------------------------------------------------------------------------

In Chapter \ref{chapter:vc}, thermal transport in two model alloys 
was investigated. The work provides several original 
insight into the physics of thermal 
transport in disordered lattices (i.e., isotopic solids and alloys). 
The first rigourous test of the virtual crystal (VC) approximation 
was presented. The VC-ALD technique has been used in a number of 
recent studies,
\cite{garg_role_2011,tian_phonon_2012,lindsay_thermal_2012,
li_thermal_2012,luckyanova_coherent_2012} 
but its limits had not been assessed until this study.  
The limits of the VC-ALD 
approach were determined using computationally-inexpensive empirical 
potentials and self-consistently treating the disorder explicitly 
and as a perturbation. 
The results indicate that while VC-ALD is generally 
an accurate method for materials 
whose thermal conductivity is dominated by low-frequency vibrational 
modes, care must be taken when modeling alloys with low thermal 
conductivities, where significant underprediction of thermal 
conductivity is likely. 

The high-scatter limit of thermal diffusivity,  
typically used in modeling amorphous materials, is directly relevant 
to the modeling of disordered lattices. 
Application of the VC 
approximation leads to vibrational mode diffusivities that are 
non-physical and the high-scatter limit provides a simple, 
physically-sound approach for correcting these predictions. 
The following calculations that were performed are novel additions 
to the literature:
\begin{itemize}
\item Using the VC-NMD method to model disordered 
lattices explicitly. This approach is limited in 
that the group velocity cannot be extracted so that thermal conductivity 
cannot be predicted.  The novel contribution is the use of 
NMD to predict the lifetimes of a disordered lattice using VC-NMD, 
where the normal modes of the Virtual Crystal (VC) are used as an 
approximation (see Section , Fig. and Appendix ).
\item  To model the disorder 
explicitly, the AF Theory calcualtions were perfomred 
on a disodered lattice. This theory has only previously been applied 
to amorphous phases.
\cite{feldman_thermal_1993,feldman_numerical_1999,
shenogin_,he_heat_2011}
The AF theory predictions showed that the lower-limit 
of diffusivity of high-frequency modes in a disordered lattice is the 
high-scatter limit, in contrast to the VC-ALD method, which incorrectly 
predicts that the limiting value is zero. Identification of this 
high-scatter limit of mode diffusivity was essential for identifying 
the breakdown in the VC-ALD method.  The high-scatter limit of 
diffusivity is usually assumed, without theoretical justification, in 
phenomological models for disordered and amorphous materials.
\cite{kittel_interpretation_1949,graebner_phonon_1986,
cahill_lattice_1988} 
This study gives self-consistent justification for its use.
\item Structure Factor of Disordered Lattice to Predict Group Velocities. 
The structure factor was calculated for modes in a model disordered 
lattice, which has previously been calculated for modes in amorphous 
materials.
\cite{biswas_vibrational_1988,feldman_thermal_1993,
allen_diffusons_1999,feldman_numerical_1999,
taraskin_determination_1999,taraskin_propagation_2000,
volz_molecular-dynamics_2000,
gotze_evolution_2000,horbach_high_2001,
martin-mayor_dynamical_2001,feldman_calculations_2002,
ciliberti_brillouin_2003,christie_vibrational_2007,
shintani_universal_2008,wyart_scaling_2010,
beltukov_ioffe-regel_2013,larkin_predicting_2013,
marruzzo_heterogeneous_2013} 
The structure factor predictions help us to 
understand that the VC-predicted group velocities are an 
underprediction of the representative velocity scale for mode 
diffusivities in the disordered lattice.  While previous studies 
have attempted to predict the group velocity of modes in disordered 
systems, there is no theoretical justification for the methods used.
\cite{duda_reducing_2011,donadio_atomistic_2009,
he_heat_2011,he_thermal_2011,he_morphology_2011,hori_phonon_2013} 
The structure factor provides a rigorous manner to estimate 
group velocities and is a significant contribution to understanding 
how to predict the correct velocity scale for mode diffusivities in 
disordered systems. 
\end{itemize}

By using all four methods discussed in 
Section \ref{S-SystemMode}, a self-consistent study of the VC approximation 
identified important connections between the modeling 
of disordered lattices and amorphous materials. 
By using three phases of LJ argon (perfect crystal, disordered lattice, 
and amorphous phase), the applicability 
of the different methods for predicting the thermal conductivity and 
mode-properties was demonstrated:
\begin{itemize}
\item MD-based GK method: suitable for modeling all 
three phases, but does not predict the mode properties.
\item Phonon based VC-ALD and VC-NMD: suitable for the perfect crystal 
and disordered lattices with the high-scatter limit correction.
\item The AF theory of diffusons: suitable for the high-frequency modes 
of the disordered lattice and all modes of the amorphous phase. 
We are unaware of any other study that uses all five of these methods 
self-consistently on the same material system. Our work provides clear 
guidelines for others on what tools are appropriate for different 
solid state systems.
Contribution to the Literature
Due to their low thermal conductivities, alloys are currently an active 
area of research, notably in the thermoelectric energy conversion field. 
The ability to predict alloy thermal conductivity is critical in 
narrowing down a large materials design space. Recent papers 
[e.g., PRL 106, 045901 (2011), PRL 109, 095901 (2012), 
PRB 85, 184303 (2012)] have used the VC-ALD method to make such 
predictions.
We believe this work will make an important contribution to the 
literature because the high-scatter limit adjustment is of interest 
to the study of low-thermal conductivity alloys. Thermoelectric energy 
generation materials, such as PbTe/Se alloys [21,22,49], maximize their 
efficiency by minimizing their thermal conductivity. The search for 
lower thermal conductivity alloys will require the modeling of even 
lower-thermal conductivity alloys, where the high-scatter limit we 
have proposed should be considered.
\end{itemize}

% The breakdown of the VC-ALD method has gone unnoticed in previous 
% computational work because these studies: 
% \begin{itemize}
% \item were limited to the 
% VC-ALD method because of computationally-expensive DFT calculations, 
% so that validation was not possible [12,18-27].
% \item focused on 
% materials where the thermal conductivity is dominated by 
% low-frequency vibration modes [12,18-27]
% \item did not always compare their 
% predictions with experimental measurements.\cite{tian_phonon_2012}
% \end{itemize}

% In our work, we provide a self-consistent study of thermal transport 
% in disordered lattices using a set of complementary computational tools 
% based in molecular dynamics simulations and lattice dynamics calculations. 
% The use of empirical potentials versus computationally-expensive DFT 
% calculations allowed us to perform the molecular dynamics simulations 
% that were necessary to observe the breakdown of the VC-ALD method. Our 
% study includes two test materials that demonstrate the applicability 
% and breakdown of the VC-ALD method. The conclusions are of general use 
% for the study of any disordered lattice. 

%--------------------------------------------------------------------------
\subsection{\label{Overview:MFP}
Mean Free Paths of Propagating Modes in Disordered Materials}
%--------------------------------------------------------------------------

A clear theoretical and modeling framework for amorphous materials was 
presented, which can form the basis for studying any disordered material. 
Drew on results and theory from the amorphous literature, which 
is extensive. 
Much work has been done by the physics community on vibrational spectra 
in amorphous solids, particularly on a-SiO$_2$. However, few of these 
studies investigated thermal transport properties. Only recently 
have thermal transport properties been studied.\cite{}

The mode lifetimes can be predicted using NMD.
If the model is large and low enough frequencies are accessed, the sound 
speed can be used. This can establish an upper bound on the mode 
diffusivities. 

The NMD lifetimes can be used with the AF theory diffusivities to predict 
effective diffuson MFPs and group velocities, which are shown in insets 
of Figs. \ref{F:vAF}(a) and \ref{F:vAF}(b). As noted in Section , the 
diffuson MFPs range between the system size $L$ and the lattice constant 
$a$ (see Section ). 
For diffusons, this effective diffuson velocities range from the 
sound speed to a fraction of it (about 0.3 for a-Si, see Fig. ). These 
diffuson velocities are much larger than 
those predicted in previous studies where the effective group velocities 
for some modes were approximately zero.\cite{}

\begin{itemize}
\item Identified the effects of metastability in amorphous materials 
on predicting lifetimes using the NMD method 
(see Section and Appendix ). 
Metastability is likely 
to affect the application of the NMD method in other ordered and 
disordered systems with weak bonding (see Section \ref{Future:LUC}). 
\item Identified differnces in the structural properties of 
of a-SiO$_2$ and a-Si that lead to a substantial difference in the 
propagating contribution to thermal conductivity in each. 
\item Predicted effective dispersion from the static structure factor (see 
Section ). The effective dispersion, when combined with diffusivity 
predictions from the AF theory and NMD method, are able to predict 
self-consistently the mode group velocities in a disorded system. 
Previous attempts to predict mode group velocities lack a theoretical basis. 
Presented a clear theoretical and computational framework 
for predicting the mode-level properties in amorphous materials. 
By predicting the effective dispersion using the static structrue 
factor, showed that the sound speed is the valid for 
low frequency modes. By using the NMD and AF-predicted 
diffusivities, demonstrated that the effective group velocity of 
diffuson modes is larger than predictions from recent studies.
(cite) 
\item Demonstrated that the NMD-predicted diffusivities are more 
reliable than those predicted by the AF theory at low frequencies 
(see Section ).
\item By comparing predictions from the models of a-Si using 
both $\omega^{-2}$ and $\omega^{-4}$ scalings, showed that 
further experimentation for a-Si thin films.

\end{itemize}


%--------------------------------------------------------------------------
\begin{figure}
\begin{center}
\centering
\includegraphics[scale=1.0]
{/home/jason/disorder/si/amor/m_af_si_normand_4096_Lambda_3.eps}
\end{center}
\caption{\label{F:vAF} 
MFPs of propagons and diffusons in a-SiO$_2$ and a-Si (see Section ).
}
\end{figure}
%--------------------------------------------------------------------------

\clearpage

%--------------------------------------------------------------------------
\subsection{\label{S:Predictive}
Predictive ability versus Computational Cost (WORK) (READ)}
%--------------------------------------------------------------------------

With the results from all studies...

now prepared to make a ranking of predictive capability. 

The GK method played an imporant role in verifying the mode properties 
prediced by all methods. In Section, the GK method provided a common 
comparison for the $\Phi$ and $\Phi'$ methods, which helps to confirm 
the disagreement between the two methods. In Section, the GK method 
provides a comparison for predictions from the VC-NMD and VC-ALD 
methods, which helps to identify the high-scatter limit of 
the diffuson mode diffusivities in disordered lattices.(cite) 
Finally, the GK method helps to confirm the scaling of the 
low-frequency contribution of the finite models of a-Si which 
is used to extrapolate to a bulk (infinte)-size value (Section ). 
The GK method will most likely be a vaulable modeling tool 
for future work. 

The VC-ALD method was shown to be limited to low frequency modes 
and best suited to high-thermal conductivity materials. High thermal 
conductivity (low-frequency dominated) materials are well-modeled 
by VC-ALD. VC-ALD is not well-suited to low thermal conductivity 
(full spectrum) materials. 
The AF theory models accurately the mid- and high-frequency 
modes in disordered materials, it does not handle the low-frequency 
modes for disordered lattices. The predictions for amorphous 
materials is also inconsistent. 
The VC-ALD method and AF theory can be supplimented by NMD calcualtions. 
The VC-NMD method is able to accurately predict the lifetimes 
of all vibrons in disordered lattices (see Section ). Propagating modes 
in a-Si can be identified definitively by NMD-predicted 
lifetimes (see Section ), while the AF theory does not show a 
definitive scaling (see Section ). 

With these findings, the theoretical methods (first discussed in 
Section ) are ranked in order of their predictive capabilites are in 
Table \ref{T-comparison-predictive}. The NMD, while the most 
computationally demanding of the four theoretical methods 
(Table \ref{T-comparison}), 
is ranked first in mode- and second in system-level predictive 
capability. The reasons for these rankings are:

\begin{itemize}
\item The NMD method is derived correctly from anharmonic lattice dynamics 
theroy (see Section ).. NMD predicts accurately the mode lifetimes and 
theraml conductivities compared to the $\Phi'$ method (see Section ).
\item The VC-NMD method accurately predicts the mode lifetimes for 
disordered lattices compared to the VC-ALD method. This leads to better 
agreement with the GK method, which is the most accurate system-level 
method.
\item The NMD method accurately predicts the low-frequncy lifetimes 
for a-Si, while the AF theory predictions have large fluctuations and 
depend on the broadening factor. The scaling from the NMD lifetimes is 
used to extrapolate a bulk thermal conductivity which is in good 
agreement with the system-level GK method (see Section ). 
\end{itemize}

The ALD and AF theory are considered to be equivalent at predicting 
mode-level and system-level properties because, for disordered lattices 
(alloys, see Section ), the VC-ALD method fails to accurately 
predict the mode lifetimes at high frequencies, 
while the AF theory is invalid at low frequencies 
(see Section ). Either 
method could be considered superior depending on whether the material 
being studied is low-frequency dominated or full-spectrum 
(see Section ). 

%--------------------------------------------------------------------------
\begin{center}
%\squeezetable
\begin{table}
\small
\caption{\label{T-comparison-predictive}Ranking of the predictive ability of 
theoretical techniques for mode-level and system-level thermal 
properties for disordered systems.}
%\begin{tabular}{llllll}
\begin{tabular}{p{0.9in}|p{0.9in}|p{0.9in}|p{0.9in}|p{0.9in}|p{0.9in}}
\hline\hline
System-level
&CH Theory 
&ALD 
&AF Theory 
&NMD 
&GK \\ 
\hline
Mode-level
&CH Theory 
&GK
&AF Theory
&ALD
&NMD \\ 
\hline\hline
\end{tabular}
\end{table}
\end{center}
%--------------------------------------------------------------------------

%--------------------------------------------------------------------------
\section{\label{Future}Future Work}
%--------------------------------------------------------------------------

%--------------------------------------------------------------------------
\subsection{\label{Future:LUC}Large Unit Cell Materials (WORK) (READ)}
%--------------------------------------------------------------------------

Large unit cell (LUC) materials are an important class of crystalline 
materials with a wide range of transport applications.(cite) 
Large unit cell crystals have an ordered (crystalline) structure, but 
the basic building block (unit cell) of the crystal has a large number 
of distinct atoms. 
They are effectively disordered 
over length scales on the order of the atomic spacing and their thermal 
conductivities can be as low as a glass.\cite{cahill_lower_1992} 
The key advantage of LUC materials is that they are still ordered 
from the standpoint of electrons, which results in good thermoelectric 
performance.
\cite{he_thermoelectric_2006,yang_effect_2006,wang_thermoelectric_2007} 

Some LUC 
materials, such as SiO$_2$-based zeolites, have been 
long-studied. Others, such as C$_60$(cite) or 
\href{http://en.wikipedia.org/wiki/Phenyl-C61-butyric_acid_methyl_ester}
{PCBM}, are currently being investigated for their thermal properties.
\cite{duda_exceptionally_2013} 
While experimental measurements of PCBM demonstrates that propagating 
modes contribute negligibly, the mechanism for its 
$\emph{exceptionally-low}$ is still not understood. Modeling could 
provide the necessary insights. 

From a modeling perspective, LUC materials pose a number of challenges,
theoretically and computationally, as compared to small unit cell 
(SUC) materials:

\begin{itemize}
\item LUC are crystalline, but are often organic or 
organic/inorganic hybrid materials. The structure of LUC 
materials is often poly- or quasi-crystalline, with less 
long-range order than SUC materials.(cite) 

\item Predicting model-level properties  
using ALD is challenging because the computational time 
scales as $b^4$ (see Section \ref{S-CompCost}). 

\item The presence of weak bonding in organic/inorganic 
materials can lead to metastability 
(see Appendix \ref{A:metastability}), which  
makes it challenging to perform the NMD method. 

\item MD simulations of LUC materials also face challenges. 
Based on the classical MD simulations of disordered 
materials studied in this work, systems sizes of 100-1000 atoms 
and simualtion times of up to 1 ns are expected to be 
necessary to accurately predict thermal conductivity. DFT calcualtions 
are too computationally expensive for such MD simulations.(cite)
Because many LUC materials have complex bonding environments, 
even the classical interatomic 
potentials are often computationally-expensive.(cite)
\end{itemize}

Based on the results in this work, there are several 
modeling strategies that can be used to study LUC materials:

\begin{itemize}
\item Identify signs of propagating modes from experimental 
measurements, if available.(cite) 
\item Based on the results for alloys (see Section ), the AF 
diffuson theory may have application for LUC materials, particularly 
at high frequencies and for 
those LUC materials which are only quasi-crystalline, 
such as PCBM.(cite)
\item The high-scatter limit for thermal conductivity (Eq. ) 
can be used to establigh a plasuible lower-bound for LUC 
materials.(cite) Similarly, the 
high-scatter limit for mode diffusivity can establish 
lower-bounds on the mode-level properties. 
\item Perform calcualtions using computationally-inexpensive 
classical interatomic potentials to asses if DFT calculations 
are necessary. 
\end{itemize}


%--------------------------------------------------------------------------
\begin{figure}
\begin{center}
\centering
\includegraphics[scale=0.4]
{/home/jason/thesis/thesis/duda_exceptionally_2013_fig4.eps}
\end{center}
\caption{\label{F:supercells} 
Range of thermal conductivity for crystalline and amorphous materials. 
Recent measurements on LUC materials show thermal conductivities 
much lower than bulk amorphous materials. 
Reproduced from Ref. \citenum{duda_exceptionally_2013}.
}
\end{figure}
%--------------------------------------------------------------------------

\clearpage

%--------------------------------------------------------------------------
\subsection{\label{Future:Timescales}
Timescales from Larger MD Simulations (WORK) (READ)}
%--------------------------------------------------------------------------

%--------------------------------------------------------------------------
\subsubsection{\label{Future:Timescales:NMD}
Exact Normal Modes}
%--------------------------------------------------------------------------

The NMD method used throughout this work is limited by the computational 
demands, which require a larger number ($\sim$ 100) of parallel processors 
to perform the analysis in a reasonable amount of time 
(less than 24 hours, see Appendix \ref{A:coding_lang:Execute:Scaling}). 
While the NMD method is trivially-parallelizable over the normal 
modes, the eigenvalue solution of the normal modes themselves 
is more computationally demanding. 
Standard routines for eigenvalue solutions of the dynamical matrix 
(see Section ) can 
calculate the exact normal modes for systems up to 8000 atoms 
in less than 24 hours using current computational resources.
\cite{gale_general_2003}  
These eigenvalue solution routines typically run on single 
processors. 
The eigenvalue solutions can be performed using the 
\href{http://www.mcs.anl.gov/petsc/}{Portable, Extensible 
Toolkit for Scientific Computation}(PETSc), which has routines for 
performing eigenvalue solutions in parallel. The PETSc pacakge has 
\href{http://wiki.python.org/moin/IntegratingPythonWithOtherLanguages}
{Python bindings} contained in the 
\href{https://code.google.com/p/petsc4py/}{petsc4py} 
pacakage, which allows for easy interface with existing MD packages
such as LAMMPS.\cite{plimpton_fast_1995}  
Here is a tutorial on 
\href{https://github.com/thehackerwithin/PyTrieste/wiki/F2Py}
{using the f2py package}, and a general tutorial on 
\href{http://docs.scipy.org/doc/numpy/user/c-info.python-as-glue.html}
{using Python as a ``glue'' language}. 

The controlling Python script would work as follows...

\begin{itemize}
\item The GULP package could be 
compiled as a Python module using the 
\href{https://code.google.com/p/f2py/}{f2py} package. 
\item The dynamical 
matrix could then be transferred to the parallel eigenvalue solution 
routines contained in PETSc using petsc4py.
\item The atomic trajectories are generated using the LAMMPS 
package with the liblammps Python module.
\item The 
\end{itemize}

A previous studied investigated the vibrational normal modes 
which has been used in a previous study on LJ argon to compute a 
(cite) 
small subset of the normal modes for a system of 32,000 atoms
\cite{mazzacurati_low-frequency_1996} 
using the 
\href{http://en.wikipedia.org/wiki/Lanczos_algorithm}
{Lanczos algorithm}.\cite{golub_matrix_2012} 
The same 32,000 system was subsquently studied using MD simulations 
and the dynamic structure factor,\cite{ruocco_relaxation_2000} 
which was discussed in Section .

While parallel eigenvalue solution can increase the system sizes 
accesible with NMD, the method is ultimately limited by the 
poor scaling of the run time and memory requirements (see Section ). 

%--------------------------------------------------------------------------
\subsubsection{\label{Future:Timescales:DSF}
Dynamic Structure Factor}
%--------------------------------------------------------------------------

MD simulations are computationally efficient.  Systems sizes of nearly 
$10^6$ atoms have been studied in this work, which were bulk systems 
with equal simulation side lengths in all three spatial dimensions. 
The dynamic structure factor, described in Section \ref{Appendix_B}, 
can predict vibrational timescales and does not 
need the eigenvectors of the exact normal modes to perform the mapping 
of the atomic trajectories. This method, 
combined with appropriately shaped supercells, could probe the 
timescales of vibrational modes up to wavelengths between 
24 and 100 nm using similar computational resources to those 
used in this work (see Section ). 
This presents an oppurtunity to compare with recent experimental 
measurements spectral linewidths in the frequency range below 
1 THz, which have been recently reported for a-SiO$_2$ but are 
lacking for a-Si.\cite{hondongwa_ultrasonic_2011} 
The current correlation function, closely related to the dynamic 
structure factor,\cite{horbach_high_2001} can even be used to study 
the spectral character of motions in a fluid.\cite{boom_molecular_1980}  

%--------------------------------------------------------------------------
\subsection{\label{Future:DFT}DFT and Thermal Modeling}
%--------------------------------------------------------------------------

Try MD and AF using better scaling DFT packages such as CP2K and BigDFT. 

BigDFT is linear scaling, but not until a large system size is reached. 

%--------------------------------------------------------------------------
\subsection{\label{Future:Package}
Comprehensive Package for Thermal Transport Calculations (WORK) (READ)}
%--------------------------------------------------------------------------

Four different methods were used in this work. No single package 
exists to perform the necessary calcaultions for all these methods. 


There exist numerous packages for performing calculations necessary for 
predicting thermal transport properties. No one package can perform all 
calculations necessary, particularly both the 
mode-level (bottom-up) and system-level (top-down) thermal transport 
properties.   

\begin{center}
\begingroup
%\squeezetable
\begin{table}
\caption{\label{T:available_codes}
Available packages
}
%\begin{ruledtabular}
\begin{tabular}{llllll}
\hline
Package & Language & Capabilities \\
\hline
GULP & Fortran & mode-level (harmonic) \\
LAMMPS & C++ & system-level, mode-level (limited) \\
phonopy & Python & mode-level (harmonic) \\
\end{tabular}
%\end{ruledtabular}
\end{table}
\endgroup
\end{center}

LAMMPS, for example, contains most methods for predicting the 
system-level thermal conductivity. A package to predict the 
mode-level properties is needed.  Ideally, the mode-level properties 
could be predicted alongside the system-level calculations, as is 
necessary to perform the NMD and spectral techniques described in 
Sections and . 

The Python language is an ideal environment for ''gluing`` together 
the availble codes, and extending their functionality in dynamic ways. 

LAMMPS as a model of social interaction, mailing list. 


num-gulp-citations(year) = [1996 1997 1998 1999 2000 2001 2002 2003 2004 2005 2006 2007 2008 2009 2010 2011 2012 2013] 

num-gulp-citations-year = [1 1 1 1 15 65 59 68 88 77 92 98 114 116 139 128 139 79*(12/9)] 

